\section{Assumptions and Scope}

The theorem uses a fixed rectangular Cartesian grid with periodic indexing in both coordinate directions. The hydrodynamic bed is flat and fixed. No source term from bathymetry is present, and no well-balanced claim is made.

The water depth is assumed strictly positive at the beginning of the step:
\[
  m^n = \min_{i,j} h_{i,j}^n > 0.
\]
The gravitational constant satisfies \(g>0\), and the mesh and time-step parameters satisfy
\[
  \Delta x>0,\qquad \Delta y>0,\qquad \Delta t\ge 0 .
\]
The suspended sediment satisfies \(s_{i,j}^n\ge 0\), the settling velocity satisfies \(w_s\ge 0\), and the bed porosity parameter satisfies \(0\le \lambda_p<1\).

The model is dilute: sediment does not alter water density, pressure, or the shallow-water wave speeds. The suspended sediment is transported by the post-hydrodynamic velocity field and then deposited exactly over one time step. The diffusivity is set to zero in the theorem:
\[
  D_s = 0.
\]

The following features are excluded from the theorem and from the numerical diagnostics reported here:
\[
\begin{gathered}
\text{nonflat bathymetry, well-balancedness, diffusion, open or closed boundaries,}\\
\text{moving-bed hydrodynamic feedback, wetting/drying, tides, erosion, resuspension,}\\
\text{higher-order reconstruction, and physical validation.}
\end{gathered}
\]
These restrictions are not incidental. They are what allow the first proof to be stated without hidden boundary constructions or unproved coupling assumptions.
